\documentclass[11pt,a4paper]{article}

% ---------- Encodage et langue ----------
\usepackage[utf8]{inputenc}
\usepackage[T1]{fontenc}
% \usepackage{lmodern}  % décommenter si disponible (police vectorielle)
\usepackage[french]{babel}
\usepackage[a4paper,margin=2.3cm]{geometry}

% ---------- Mathématiques ----------
\usepackage{amsmath,amssymb,amsthm}
\usepackage{mathtools}
\usepackage{enumitem}
\usepackage[most]{tcolorbox}
\usepackage[colorlinks=true,linkcolor=blue!60!black,urlcolor=blue!60!black]{hyperref}

% ---------- Macros ----------
\newcommand{\R}{\mathbb{R}}
\newcommand{\N}{\mathbb{N}}
\newcommand{\Z}{\mathbb{Z}}
\newcommand{\Q}{\mathbb{Q}}
\newcommand{\C}{\mathbb{C}}
\newcommand{\Rb}{\overline{\R}}
\newcommand{\interieur}[1]{\mathring{#1}}
\newcommand{\adh}[1]{\overline{#1}}
\newcommand{\Int}{\operatorname{int}}
\newcommand{\diam}{\operatorname{diam}}
\newcommand{\dist}{\operatorname{dist}}
\newcommand{\Bf}{B_f}
\newcommand{\norm}[1]{\left\|#1\right\|}
\newcommand{\abs}[1]{\left|#1\right|}
\DeclareMathOperator{\Ker}{Ker}
\newcommand{\ie}{c'est-à-dire}
\newcommand{\K}{\mathbb{K}}
\newcommand{\Vect}{\operatorname{Vect}}
\newcommand{\opnorm}[1]{|\!|\!|#1|\!|\!|}
\newcommand{\ps}[2]{\left\langle #1,#2\right\rangle}
\newcommand{\Lc}{\mathcal{L}}
\newcommand{\Cz}{\mathcal{C}^0}
\newcommand{\Cinf}{\mathcal{C}^\infty}
\newcommand{\Cc}{\mathcal{C}^0_c}

% ---------- Environnements ----------
\newtcolorbox{exo}[1]{
  enhanced, breakable,
  colback=black!3, colframe=black!70,
  fonttitle=\bfseries, title={#1},
  boxrule=0.6pt, arc=1pt, left=6pt, right=6pt
}
\newtcolorbox{pointcle}{
  enhanced, breakable,
  colback=blue!4, colframe=blue!55!black,
  fonttitle=\bfseries, title={Point clé},
  boxrule=0.8pt, arc=2pt, left=6pt, right=6pt
}
\newtcolorbox{piege}{
  enhanced, breakable,
  colback=red!4, colframe=red!60!black,
  fonttitle=\bfseries, title={Piège / difficulté},
  boxrule=0.8pt, arc=2pt, left=6pt, right=6pt
}
\newtcolorbox{commentaire}{
  enhanced, breakable,
  colback=gray!8, colframe=gray!55!black,
  fonttitle=\bfseries, title={Commentaire},
  boxrule=0.6pt, arc=2pt, left=6pt, right=6pt
}
\newtcolorbox{methode}{
  enhanced, breakable,
  colback=green!4, colframe=green!40!black,
  fonttitle=\bfseries, title={Méthode},
  boxrule=0.8pt, arc=2pt, left=6pt, right=6pt
}

\theoremstyle{definition}
\newtheorem*{solution}{Solution}

\title{\textbf{4MA305 -- Bases d'analyse fonctionnelle}\\[2mm]
\Large TD n\textsuperscript{o}1 -- Topologie des espaces métriques (énoncés)}
\author{}
\date{Année universitaire 2026--2027}

\begin{document}
\maketitle

\begin{commentaire}
Feuille d'exercices sur le chapitre~1 du polycopié. La \textbf{partie A} (Exercices~1 à~6) est complémentaire des exercices du poly (corrigés dans le TD0). La \textbf{partie B} reproduit la \textbf{feuille officielle} du TD n\textsuperscript{o}1 (Exercices officiels~1 à~13). La \textbf{partie C} (Exercices~7 et~8) prépare l'examen. Le corrigé détaillé est dans \texttt{corrige\_TD1.pdf} ; l'idée est de chercher les exercices en conditions d'examen avant de le lire. Difficulté indicative, partie A : 1, 2, 3 (fondamentaux) ; 4, 5 (classiques d'examen) ; 6 (plus astucieux). Partie B : 1 à 6 (cours) ; 7, 8, 10, 11 (à savoir refaire) ; 13 (le plus long).
\end{commentaire}

\section*{Partie A -- Exercices complémentaires}

\begin{exo}{Exercice 1 -- Distances bornées : ce qu'une distance retient de plus que sa topologie}
Soit $(X,d)$ un espace métrique. On pose $\delta:=\min(1,d)$ et $\delta':=\dfrac d{1+d}$.
\begin{enumerate}[label=(\alph*)]
  \item Soit $\varphi:\R_+\to\R_+$ croissante, sous-additive ($\varphi(s+t)\le\varphi(s)+\varphi(t)$) et ne s'annulant qu'en $0$. Montrer que $\varphi\circ d$ est une distance sur $X$. En déduire que $\delta$ et $\delta'$ sont des distances.
  \item Montrer que l'identité de $(X,d)$ dans $(X,\delta)$ est uniformément continue ainsi que sa réciproque ; même question pour $\delta'$. En déduire que $d$, $\delta$ et $\delta'$ ont les mêmes ouverts, les mêmes suites de Cauchy, et que $(X,d)$ est complet si et seulement si $(X,\delta)$ l'est.
  \item Montrer que $X$ est borné pour $\delta$. À quelle condition existe-t-il $c>0$ tel que $c\,d\le\delta$ ?
  \item Montrer que $(X,d)$ et $(X,\delta)$ ont les mêmes parties précompactes. Que dire de $(\R,\min(1,|\cdot|))$ : borné ? complet ? précompact ? compact ?
\end{enumerate}
\end{exo}

\begin{exo}{Exercice 2 -- Séparabilité}
Un espace métrique est \emph{séparable} s'il contient une partie dénombrable dense.
\begin{enumerate}[label=(\alph*)]
  \item Soit $X$ un espace métrique séparable et $A\subset X$. Montrer que $A$, muni de la distance induite, est séparable.
  \item Montrer qu'un espace métrique précompact est séparable. En déduire que tout espace métrique compact est séparable.
  \item Soit $X$ un espace métrique séparable. Montrer que toute famille d'ouverts non vides deux à deux disjoints de $X$ est au plus dénombrable.
  \item Soit $X:=\{0,1\}^\N$ muni de $d_\infty(u,v):=\sup_n|u_n-v_n|$. Montrer que $X$ n'est pas séparable. En déduire que $\ell^\infty(\N)$ ne l'est pas.
\end{enumerate}
\end{exo}

\begin{exo}{Exercice 3 -- Uniforme continuité : critère séquentiel, exemples, prolongement}
\begin{enumerate}[label=(\alph*)]
  \item Soient $(X,d)$, $(Y,d')$ deux espaces métriques et $f:X\to Y$. Montrer que $f$ est uniformément continue si et seulement si, pour toutes suites $(x_n)$, $(y_n)$ de $X$ telles que $d(x_n,y_n)\to0$, on a $d'(f(x_n),f(y_n))\to0$.
  \item Les fonctions $x\mapsto x^2$ et $x\mapsto\sin(x^2)$ sont-elles uniformément continues sur $\R$ ? La fonction $x\mapsto\sqrt x$ est-elle uniformément continue sur $\R_+$ ? lipschitzienne ?
  \item Soit $f:\,]0,1]\to\R$ continue. Montrer que $f$ est uniformément continue si et seulement si elle admet une limite finie en $0$. Application : $\sin\frac1x$ et $x\sin\frac1x$.
  \item Soit $f:\R\to\R$ continue admettant des limites finies en $+\infty$ et en $-\infty$. Montrer que $f$ est uniformément continue.
\end{enumerate}
\end{exo}

\begin{exo}{Exercice 4 -- Trois applications du lemme de Baire}
\begin{enumerate}[label=(\alph*)]
  \item Soit $X$ un espace métrique complet, non vide et sans point isolé. Montrer que $X$ n'est pas dénombrable. En déduire qu'il n'existe aucune distance sur $\Q$ définissant sa topologie usuelle et le rendant complet.
  \item Soit $X$ un espace métrique complet non vide et $(f_n)$ une suite de fonctions continues de $X$ dans $\R$ telle que, pour tout $x\in X$, la suite $(f_n(x))_n$ soit bornée. Montrer qu'il existe un ouvert non vide $U\subset X$ et $M\ge0$ tels que $|f_n(x)|\le M$ pour tout $n$ et tout $x\in U$.
  \item Soit $f:\,]0,+\infty[\,\to\R$ continue telle que, pour tout $x>0$, $f(nx)\to0$ quand l'entier $n$ tend vers l'infini. Montrer que $f(x)\to0$ quand $x\to+\infty$.
\end{enumerate}
\end{exo}

\begin{exo}{Exercice 5 -- Théorème de Dini}
Soit $X$ un espace métrique compact et $(f_n)$ une suite \emph{croissante} de fonctions continues de $X$ dans $\R$, convergeant simplement vers une fonction $f$ \emph{continue}.
\begin{enumerate}[label=(\alph*)]
  \item Montrer que la convergence est uniforme. \emph{Indication : pour $\varepsilon>0$, considérer les fermés $K_n:=\{x:f(x)-f_n(x)\ge\varepsilon\}$.}
  \item Montrer par des exemples que chacune des trois hypothèses (compacité de $X$, continuité de la limite, monotonie) est nécessaire.
  \item On définit des polynômes par $P_0:=0$ et $P_{n+1}(x):=P_n(x)+\frac12\big(x-P_n(x)^2\big)$. Montrer que $0\le P_n\le P_{n+1}\le\sqrt x$ sur $[0,1]$, puis que $(P_n)$ converge uniformément vers $\sqrt x$ sur $[0,1]$. En déduire que $t\mapsto|t|$ est limite uniforme de polynômes sur $[-1,1]$.
\end{enumerate}
\end{exo}

\begin{exo}{Exercice 6 -- Isométries et applications dilatantes d'un compact}
Soit $(X,d)$ un espace métrique compact non vide et $f:X\to X$.
\begin{enumerate}[label=(\alph*)]
  \item On suppose que $f$ est une isométrie : $d(f(x),f(y))=d(x,y)$ pour tous $x,y$. Montrer que $f$ est surjective. \emph{Indication : si $a\notin f(X)$, étudier la suite $(f^n(a))_n$.}
  \item Donner un exemple d'isométrie non surjective d'un espace métrique complet dans lui-même. En déduire une nouvelle preuve que la boule unité fermée de $\ell^2(\N)$ n'est pas compacte.
  \item On suppose seulement que $f$ est \emph{dilatante} : $d(f(x),f(y))\ge d(x,y)$ pour tous $x,y$. Montrer que $f$ est une isométrie. \emph{Indication : pour $x,y$ fixés, extraire une sous-suite convergente de $\big(f^n(x),f^n(y)\big)_n$ dans $X\times X$.}
  \item Soit $g:X\to X$ surjective et $1$-lipschitzienne. Montrer que $g$ est une isométrie.
\end{enumerate}
\end{exo}

\section*{Partie B -- Feuille officielle du TD n\textsuperscript{o}1}

\subsection*{Topologie générale d'un espace métrique}

\begin{exo}{Exercice officiel 1 -- Intérieur, adhérence et frontière ne commutent pas}
Soit $X$ un espace métrique et $A$ une partie de $X$. Donner des contre-exemples aux propriétés suivantes :
\begin{enumerate}[label=\arabic*.]
  \item $\interieur{\adh{A}} \subset A$,
  \item $\interieur{\adh{A}} \supset A$,
  \item $\partial A \subset \partial \adh{A}$.
\end{enumerate}
\end{exo}

\begin{exo}{Exercice officiel 2 (Exercice 1.3 du poly) -- Topologie induite}
Soit $(X,d)$ un espace métrique et $A\subset X$. Vérifier que les ouverts et fermés de $A$ (pour la distance induite) sont respectivement les intersections avec $A$ des parties ouvertes et fermées de $X$.
\end{exo}

\begin{exo}{Exercice officiel 3 -- Vrai ou faux sur la topologie induite}
Soit $X$ un espace métrique, $A\neq\emptyset$ une partie de $X$. Dire si les propriétés suivantes sont vraies ; si elles sont fausses, donner des contre-exemples et, éventuellement, des conditions supplémentaires pour qu'elles soient vraies.
\begin{enumerate}[label=\arabic*.]
  \item Si $B\subseteq A$ est un ouvert de $A$ alors $B$ est aussi ouvert dans $X$.
  \item Si $B$ est ouvert dans $X$, $B\cap A$ est aussi ouvert dans $A$.
  \item Si $B\subseteq A$ est fermé dans $A$, $B$ est aussi fermé dans $X$.
  \item Si $B\subseteq A$ est fermé dans $X$, $B$ est aussi fermé dans $A$.
  \item Si $X$ est un espace métrique complet, $A$ est complet.
  \item Il existe une fonction $f:\Q\to\N$ non constante, continue.
\end{enumerate}
\end{exo}

\begin{exo}{Exercice officiel 4 -- Boules ouvertes, boules fermées, sphères}
Soit $(X,d)$ un espace métrique, $a\in X$ et $r\ge0$. On note $B(a,r)$, $\Bf(a,r)$, $S(a,r)$ la boule ouverte, la boule fermée et la sphère de centre $a$ et de rayon $r$.
\begin{enumerate}[label=\alph*)]
  \item Montrer que $B(a,r)$ est ouvert, que $\Bf(a,r)$ est fermé et que $S(a,r)$ est fermé.
  \item Montrer que $\adh{B(a,r)}\subset\Bf(a,r)$ et $B(a,r)\subset\Int(\Bf(a,r))$. Montrer que dans un espace vectoriel normé ce sont des égalités. Donner des exemples d'espaces métriques où ces inclusions sont strictes.
\end{enumerate}
\end{exo}

\begin{exo}{Exercice officiel 5 (Exercice 1.7 du poly) -- Convergence dans un produit}
Soient $(X,d_X)$ et $(Y,d_Y)$ deux espaces métriques. On munit $X\times Y$ de la distance
\[
D_1\big((x_1,y_1),(x_2,y_2)\big):=d_X(x_1,x_2)+d_Y(y_1,y_2).
\]
Montrer que la convergence de $z_n:=(x_n,y_n)$ vers $z:=(x,y)$ équivaut aux deux convergences $x_n\to x$ dans $X$ et $y_n\to y$ dans $Y$.
\end{exo}

\begin{exo}{Exercice officiel 6 (Exercice 1.11 du poly) -- La distance à une partie est $1$-lipschitzienne}
Soit $(X,d)$ un espace métrique et $A\subset X$ une partie non vide. Montrer que $x\mapsto d(A,x):=\inf_{a\in A}d(x,a)$ est $1$-lipschitzienne.
\end{exo}

\subsection*{Complétude}

\begin{exo}{Exercice officiel 7 -- Deux applications du lemme de Baire}
\begin{enumerate}[label=\arabic*.]
  \item Soit $E$ un espace métrique complet et $(F_n)_{n\in\N}$ une suite de fermés de $E$ telle que $E=\bigcup_{n}F_n$. Montrer que $\bigcup_n\interieur{F_n}$ est un ouvert dense de $E$.
  \item Soit $T:E\to E$ linéaire continue, $E$ normé complet. On suppose que pour tout $x\in E$ il existe $n_x$ tel que $T^{n_x}x=0$. Montrer que $T$ est nilpotente.
\end{enumerate}
\end{exo}

\begin{exo}{Exercice officiel 8 -- Théorème de Cantor (Exercice 1.14 du poly)}
Montrer que pour un espace métrique $(X,d)$, la complétude équivaut à la propriété suivante : si $(F_n)_{n\in\N}$ est une suite décroissante de fermés non vides vérifiant $\lim_n\diam(F_n)=0$, alors $\bigcap_nF_n$ est un singleton.
\end{exo}

\subsection*{Compacité}

\begin{exo}{Exercice officiel 9 -- Fermés bornés non compacts}
Donner des exemples de fermés bornés d'un espace métrique qui ne sont pas compacts.
\end{exo}

\begin{exo}{Exercice officiel 10 (Exercice 1.19 du poly) -- Unique valeur d'adhérence dans un compact}
Soit $(X,d)$ un espace métrique compact et $(x_n)_n$ une suite d'éléments de $X$ ne possédant qu'une valeur d'adhérence $x$. Montrer que $(x_n)_n$ converge vers $x$.
\end{exo}

\begin{exo}{Exercice officiel 11 -- Distance entre deux compacts}
Soit $(X,d)$ un espace métrique. On définit la distance entre deux sous-ensembles $A,B\subset X$ par $\dist(A,B):=\inf_{a\in A,\,b\in B}d(a,b)$. Montrer que si $A$ et $B$ sont compacts (et non vides), alors il existe $a_*\in A$ et $b_*\in B$ tels que $\dist(A,B)=d(a_*,b_*)$.
\end{exo}

\begin{exo}{Exercice officiel 12 -- Fonction continue tendant vers $0$ à l'infini}
Soit $f:\R\to\R$ continue telle que $\lim_{x\to-\infty}f(x)=\lim_{x\to+\infty}f(x)=0$. Montrer que $f$ est bornée et atteint au moins une de ses bornes.
\end{exo}

\begin{exo}{Exercice officiel 13 -- Points fixes}
\begin{enumerate}[label=\arabic*.]
  \item Soit $E$ un espace vectoriel normé et $C\subseteq E$ un convexe compact (non vide). Montrer que toute application $1$-lipschitzienne de $C$ dans lui-même admet un point fixe. \emph{Indication : considérer, pour $z\in C$ fixé, $f_n(x):=\frac{1}{n+1}z+\frac{n}{n+1}f(x)$.}
  \item Soit $(X,d)$ un espace métrique compact (non vide) et $f:X\to X$ telle que $d(f(x),f(y))<d(x,y)$ pour $x\ne y$.
  \begin{enumerate}[label=(\alph*)]
    \item Expliquer pourquoi le théorème de point fixe de Picard ne s'applique pas.
    \item Montrer cependant que $f$ admet un unique point fixe $a\in X$.
    \item Montrer que la suite des itérées de $f$ converge simplement vers la fonction constante égale à $a$.
    \item Que dire dans le cas de l'inégalité large $d(f(x),f(y))\le d(x,y)$ ?
  \end{enumerate}
\end{enumerate}
\end{exo}

\section*{Partie C -- Vers l'examen : questions de cours, vrai ou faux}

\begin{exo}{Exercice 7 -- Questions de cours}
\emph{À traiter sans document : énoncé complet, puis démonstration.}
\begin{enumerate}[label=(\alph*)]
  \item Énoncer et démontrer le théorème de prolongement des applications uniformément continues (Théorème~1.1).
  \item Énoncer et démontrer le lemme de Baire (Lemme~1.1).
  \item Énoncer et démontrer le théorème de Heine (Théorème~1.3).
  \item Montrer que dans un espace métrique complet, une partie est compacte si et seulement si elle est fermée et précompacte (Corollaire~1.4 ; on redémontrera les Propositions~1.13 et 1.14).
  \item Énoncer et démontrer le théorème de point fixe de Picard (Exercice~1.16).
\end{enumerate}
\end{exo}

\begin{exo}{Exercice 8 -- Vrai ou faux ?}
Pour chaque affirmation, dire si elle est vraie ou fausse, et justifier par une preuve courte ou un contre-exemple. $(X,d)$ et $(Y,d')$ désignent des espaces métriques.
\begin{enumerate}[label=\arabic*.]
  \item L'adhérence de la boule ouverte $B(a,r)$ est la boule fermée $\Bf(a,r)$.
  \item Dans un espace métrique complet, toute partie fermée et bornée est compacte.
  \item L'image réciproque d'un compact par une application continue est compacte.
  \item L'image d'un fermé par une application continue est fermée.
  \item Une application continue transforme les suites de Cauchy en suites de Cauchy.
  \item Si $X$ est compact et $f:X\to Y$ est continue et bijective, alors $f^{-1}$ est continue.
  \item Si $X$ et $Y$ sont homéomorphes et $X$ est complet, alors $Y$ est complet.
  \item Dans un espace métrique complet, une intersection dénombrable d'ouverts denses est un ouvert dense.
  \item $\Q$ est une intersection dénombrable d'ouverts de $\R$.
  \item Tout espace métrique compact est séparable.
  \item Dans un espace métrique, toute partie fermée et précompacte est compacte.
  \item Si $D$ est dense dans $X$ et $f:D\to Y$ est uniformément continue, alors $f$ se prolonge en une application continue $X\to Y$.
\end{enumerate}
\end{exo}

\end{document}
